\renewcommand{\thechapterdisplay}{Primeiro dia}
\renewcommand{\thechapterrunner}{Santificação deste mês}
\chapter[Primeiro dia --- Santificação deste mês]{Santificação deste mês}
\section*{Razões para a santificação deste mês}

As origens do mês dedicado aos mortos remontam ao Antigo Testamento, ao povo de Israel. Esse povo, que sozinho possuía o verdadeiro Espírito de Deus, não se contentou em proclamar em seus livros sagrados que orar pelos mortos era uma obra santa e salvífica, mas também quis determinar a duração de tal oração. Foi estabelecido que o luto não cessaria até que cada um dos falecidos fosse lamentado por um mês. Assim, após a morte de Jacó, seus filhos lastimaram sua perda e fizeram orações por trinta dias\footnote{\textit{Na verdade Jacó foi chorado por 70 dias pelos egípcios (Gênesis 50:3). Outros líderes de Israel foram lamentados por trinta dias, como Moises e Aarão (Deuteronômio 34:8 e Números 20:29).}}.

Inspirada por uma prática tão antiga e aprovada, a piedade dos fiéis consagrou um mês inteiro ao alívio das almas do Purgatório. E como a Igreja celebra o \textit{Dia de Todos os Fiéis Defuntos} no segundo dia de novembro, este mês pareceu o mais apropriado para esta devoção. O mês consagrado às almas do Purgatório, recomendado pelos soberanos pontífices e enriquecido com favores espirituais é publicamente celebrado por um grande número de comunidades religiosas e paróquias cristãs.

Acolhe com grande alegria o despertar deste mês que tão admiravelmente responde às necessidades do coração. Ele vos lembrará das mais ternas memórias de família, das mais sagradas promessas, das mais tocantes despedidas. Ele desenvolverá a vossa compaixão pelos familiares e amigos que, devido ao seu sofrimento e estado miserável, tornam-se mais queridos.

Sim, a dignidade dessas almas desafortunadas, a severidade de seu sofrimento, sua incapacidade de salvarem-se a si mesmas, a Glória de Deus, vosso interesse pessoal, enfim, tudo clama para que as visiteis e venhais em seu auxílio, cada dia deste mês. Ó! Não é este o mês dos suspiros e das lágrimas? Não é este o mês da caridade e do reconhecimento, o mês dos vivos e dos mortos, o mês verdadeiramente libertador? Cheio de entusiasmo no início de tal mês, uma santa exclamava ao começar os exercícios do mês de novembro: \textit{Esvaziemos o Purgatório!}

De maneira menos ambiciosa, mas não menos zelosa, digamos a nós mesmos, caro leitor: \textit{Quero aliviar muitas Almas do Purgatório, durante este Mês de bênçãos que lhes é consagrado! Eu quero, eu devo, eu posso!}


\section*{Meios de santificar este mês}

Para aproveitar bem o mês dos mortos, tomai hoje as seguintes resoluções e cumpri-as fielmente. A cada dia, pela manhã, oferecei à Deus, pelas almas do Purgatório, os méritos de vosso trabalho, de vossos sofrimentos e de vossas boas obras. Dizei ao despertar: \textit{Senhor, tudo para maior gloria Vossa e pelo alívio de meus antepassados falecidos!} Separai um tempo específico durante o dia para ler atentamente a este pequeno livro. Esta leitura esclarecerá o vosso espírito e aquecerá o vosso coração. Não deixeis nunca de fazê-la. Ide algumas vezes ao cemitério depositar sobre o túmulo de todos aqueles que vos foram próximos, com vossas lágrimas e orações, uma lembrança que os console. É o lugar próprio para se rezar e chorar! Cada semana, separai um dia especial às almas do Purgatório, quarta-feira por exemplo, e assisti à santa Missa nessa intenção. No decorrer do mês, fazei pelas almas uma esmola ao pobres e uma ou várias comunhões fervorosas.

Sim, fazei assim, alma cristã e compassiva, e ao final deste mês libertador tereis enviado para a Igreja triunfante do céu um grande número dos irmãos que gemem e choram nas chamas do Purgatório. Que motivo de consolações! Que penhor de esperança! \textit{Vamos, levantai-vos!} Dizia São Bernardo. \textit{Voai em socorro das almas dos defuntos, chamai sobre elas a clemência divina com gemidos de penitência, implorai a misericórdia com suspiros de mortificação, satisfazei por elas com o sacrifício da missa, resgatai-as com orações, e lhes abrireis as portas do Paraíso.}


\section*{Exemplo }

Eis como uma uma pessoa digna de fé conta sua cura extraordinária, obtida pela intercessão das almas do Purgatório, durante o mês de novembro.

\begin{quote}
\textit{\textit{Eu estava, já há vários anos, atingida por uma doença cruel que fazia do meu corpo um esqueleto, de minha vida um martírio, e que me conduzia insensivelmente ao túmulo. Consultei vários médicos de renome, mas todos os remédios que me prescreviam, após alguns raros momentos de alívio, me deixavam ainda mais frágil e abatida. Não podendo nada obter dos recursos da medicina, deixei de lado os medicamentos e recorri às almas do Purgatório que compreendem muito bem o mistério do sofrimento. O mês de novembro, que lhes é especialmente consagrado, iria começar. Tomei a resolução de o celebrar com todo o fervor possível, e de concluí-lo com uma boa Comunhão. Os meus pais e as pessoas piedosas que conhecia uniram suas preces às minhas. Cada dia, reunidos à noite em meu quarto, aos pés de uma imagem de São José, pedíamos com confiança duas coisas: a libertação das pobres almas do Purgatório e o alívio dos meus males. Ao final da primeira semana, experimentei uma melhora sensível, que coisa admirável! No último dia do bendito mês, eu estava na igreja, à mesa santa, inebriada de alegria, de felicidade e agradecendo. Minha cura estava completa, não havia sobrado traço algum da doença que me havia torturado por tanto tempo e que, como diziam os próprios médicos, era incurável. Graças sejam dadas às santas almas do Purgatório cuja proteção se manifestou de maneira visível a meu favor!}}
\end{quote}

Que favores vamos receber também, para nossos fiéis defuntos e para nós mesmos, se praticarmos santamente as devoções deste belo mês! Tendes coragem, e confiança!

\vspace{\baselineskip}\noindent \textbf{Oremos.} Deus bom e misericordioso, dignai-vos ouvir as fervorosas orações que vos dirigimos durante este mês de bênçãos. Nós vos oferecemos cada dia, cada hora, pelo alívio e libertação dessas almas cativas que suplicam a vós e a nós, do fundo de sua prisão tenebrosa. Senhor, chamai vossos filhos, nossos irmãos, ao repouso eterno. Que a luz que nunca se apaga brilhe sobre eles! Que descansem em paz!


\begin{center}
\textit{Requiescant in pace!}
\end{center}

Rezai agora uma dezena do terço e as orações no final deste livro.
